% Figure 12.1: two-loop Compton-scattering graphs.
% Source: physical PDF p. 527 (printed p. 504).
\begingroup
\renewcommand{\thefigure}{12.\arabic{figure}}
\setcounter{figure}{0}
\begin{figure}[tbp]
\centering
\begin{tikzpicture}[
  x=0.82cm,y=0.82cm,
  electron/.style={line width=0.72pt},
  photon/.style={line width=0.66pt,decorate,
    decoration={snake,amplitude=0.95pt,segment length=4.3pt}},
  flow/.style={-{Stealth[length=4.5pt,width=3.4pt]},line width=0.58pt},
  subgraph/.style={densely dotted,line width=0.72pt},
  panel/.style={font=\small}
]
% Common Compton skeletons: vertical electron line and two external photons.
\foreach \x in {0,6.8,13.6}{
  \draw[electron] (\x,-3.0) -- (\x,3.0);
  \draw[photon] (\x-2.35,1.65) -- (\x,0.92);
  \draw[photon] (\x-2.35,-1.65) -- (\x,-0.92);
  \draw[photon] (\x,2.15) .. controls (\x+1.15,1.95) and
    (\x+1.15,-1.95) .. (\x,-2.15);
}
% (a): second broad photon loop nested inside the first.
\draw[flow] (0,-0.25) -- (0,0.45);
\draw[photon] (0,1.48) .. controls (0.78,1.30) and (0.78,-1.30) .. (0,-1.48);
% (b): electron self-energy subgraph on the internal electron line.
\draw[flow] (6.8,-0.25) -- (6.8,0.45);
\draw[photon] (6.8,0.52) .. controls (7.35,0.42) and (7.35,-0.42) .. (6.8,-0.52);
\draw[subgraph] (6.22,-0.76) rectangle (7.48,0.76);
% (c): vertex subgraph containing the lower external-photon attachment.
\draw[flow] (13.6,0.28) -- (13.6,0.84);
\draw[photon] (13.6,0.02) .. controls (14.15,-0.08) and (14.15,-1.18) ..
  (13.6,-1.30);
\draw[subgraph] (13.00,-1.72) rectangle (14.35,0.20);
\node[panel] at (0,-3.72) {(a)};
\node[panel] at (6.8,-3.72) {(b)};
\node[panel] at (13.6,-3.72) {(c)};
\end{tikzpicture}
\caption{Some two-loop graphs for Compton scattering. Here straight lines are
electrons; wavy lines are photons. The momentum space integral for diagram
(a) is convergent, while for (b) and (c) it is divergent, due to the
subintegration associated with the subgraphs surrounded by dotted lines.}
\label{fig:12.1}
\end{figure}
\endgroup
